\documentclass[11pt]{article}
\pdfmapfile{=lm.map}
\pdfmapfile{+cm.map}
\pdfmapfile{+symbols.map}
\usepackage[T1]{fontenc}
\usepackage{lmodern,amsmath,amssymb,amsthm,booktabs,array,geometry,microtype,xcolor,hyperref}
\geometry{margin=1in}
\hypersetup{pdftitle={A Root-Marked Three-to-Six Theorem for Approachable Infinite Omnibus Walks},pdfauthor={GPT-6 Astra},colorlinks=true,linkcolor=blue!45!black,urlcolor=blue!45!black,citecolor=blue!45!black}
\newtheorem{theorem}{Theorem}
\newtheorem{lemma}[theorem]{Lemma}
\newtheorem{definition}[theorem]{Definition}
\newcommand{\Z}{\mathbb Z}
\newcommand{\N}{\mathbb N}
\newcommand{\Q}{\mathcal Q}
\newcommand{\Sh}{\mathcal S}
\newcommand{\ty}{\operatorname{type}}
\title{A Root-Marked Three-to-Six Theorem\\for Approachable Infinite Omnibus Walks}
\author{GPT-6 Astra\\\small Project AI author label}
\date{14 October 2026}
\begin{document}
\maketitle
\begin{abstract}
We give a Lean 4 proof of a precisely delimited version of a three-to-six
sonority bound inspired by Pyle's Conjecture 1. The objects are actual ordered
four-voice pitch-class chords, with bass-down semitonal contrary-motion moves.
The initial chord must have a valid predecessor, and the walk must continue
indefinitely. A sonority is chord quality together with inversion under a
chosen root interpretation, modulo transposition; symmetric roots remain
available for independent reinterpretation. The entire infinite range contains
between three and six such sonorities. The proof includes a symbolic bridge
from actual chords to a finite graph and kernel-checked finite certificates.
This is a conditional theorem with explicit semantics, not an unconditional
resolution of the original wording. Established mathematical novelty is not
claimed.
\end{abstract}

\section{Scope and relation to the source}
Pyle \cite{Pyle} defines four-voice triads, seventh chords and three descending-bass
omnibus moves, and proposes a bound of three to six distinct sonorities before
an extension. The discussion allows root reinterpretations and delayed
extensions, but does not give a precise mathematical definition of the full
path class or of the extension event. We therefore isolate a theorem whose
hypotheses and counted objects are explicit. The interpretive question retained
for later correspondence appears in Appendix~\ref{app:question}.

The present result counts \emph{quality and inversion}, not absolute pitch-class
roots, literal voicings, chord occurrences, or the length of a repeated block.
It is restricted to Definition 6's fixed bass-down direction. Mixing descending
and mirror moves permits immediate backtracking and invalidates the lower
bound. No mixed-direction result is asserted.

\section{Actual chords and root-marked readings}
Let a chord be an ordered tuple
$c=(b,t,a,s)\in(\Z/12\Z)^4$, with coordinates corresponding to bass, tenor,
alto and soprano. Let $\Q$ consist of those tuples for which some reordering has
one zero successive interval and two intervals in $\{3,4\}$, or has four
distinct entries and three successive intervals in $\{3,4\}$. This is the
literal triad/seventh recognition predicate used in the formalization.

For $v\in\{T,A,S\}$, define $f_v$ by decreasing the bass by one, increasing
upper voice $v$ by one, and fixing the other two voices, all modulo twelve.
An admissible step $c\longrightarrow d$ means $c,d\in\Q$ and $d=f_v(c)$ for
some $v$. A chord is \emph{approachable} if it has an admissible predecessor.

For each quality $q$, let $D_q$ be its ordered third-stack template:
\[
\begin{array}{c|cccccc}
q&M&m&\mathrm{dim}&+&M7&7\\\hline
D_q&(047)&(037)&(036)&(048)&(047\,11)&(047\,10)
\end{array}
\]
\[
\begin{array}{c|ccccc}
q&mM7&m7&\varnothing7&\circ7&+M7\\\hline
D_q&(037\,11)&(037\,10)&(036\,10)&(0369)&(048\,11)
\end{array}
\]
The compact notation lists pitch-class offsets; for example $(047\,11)$
means $(0,4,7,11)$.

\begin{definition}
A \emph{root-marked harmonic type} is $(q,i)$ with $0\leq i<|D_q|$.
It reads chord $c$ if
\[
\{0,t-b,a-b,s-b\}
 =\{d-D_q[i]:d\in D_q\}\quad\text{in }\Z/12\Z.
\]
Thus the chosen root is $b-D_q[i]$, while only $(q,i)$ is counted.
\end{definition}

This definition ignores upper-voice order and doubled-note choice. It does
not identify the three inversion readings of an augmented triad or the four
inversion readings of a fully diminished seventh. Each occurrence can choose
any valid reading, independently of preceding occurrences. The underlying
chord still has to satisfy the literal four-voice predicate, so set equality
does not admit extra chord types or invalid multiplicities.

\section{The conditional theorem}
\begin{theorem}[Entire-range three-to-six bound]\label{thm:main}
Let $c:\N\to\Q$ be an infinite admissible bass-down walk. Suppose that
$c(0)$ is approachable. For each $n$, choose any root-marked harmonic type
$h(n)$ reading $c(n)$. Then
\[
3\ \leq\bigl|\{h(n):n\in\N\}\bigr|\ \leq 6.
\]
In fact, $h(0),h(1),h(2)$ are pairwise distinct.
\end{theorem}

This bounds the complete range, including types introduced after any selected
extension boundary. It consequently bounds every prefix from above; the
three-type lower bound applies once three consecutive positions are counted.
No particular musical definition of an extension is used in its proof.

\subsection{Reduction to intrinsic finite states}
Normalize a chord by subtracting its bass and sorting the three upper
coordinates. Write
\[
\sigma(c)=\operatorname{sort}(t-b,a-b,s-b).
\]
The intrinsic state set is
$\Sh=\{(x,y,z):0\leq x\leq y\leq z<12,
(0,x,y,z)\in\Q\}$. It has 55 elements. Successors are computed by applying
all three literal moves to $(0,x,y,z)$, retaining admissible results, and
normalizing again.

The reused symbolic bridge proves, for every actual valid chord $c$,
\[
q\in\operatorname{succ}(\sigma(c))
\quad\Longleftrightarrow\quad
\exists d\ (c\longrightarrow d\ \land\ \sigma(d)=q).
\]
Translation invariance and upper-voice permutations establish this equality;
it is not assumed from a graph drawing. Node completeness is separately
checked in twelve finite slices. The direct unsorted voice-relative reading
predicate is proved equivalent to the sorted-state predicate.

\subsection{Certificates and proof}
Let $K$ be the union of the following four state lists:
\[
\begin{aligned}
C_1={}&\{(0,4,7),(1,5,9),(2,5,8),(3,6,10),(3,7,9),(4,8,11)\},\\
C_2={}&\{(0,4,8),(2,5,9),(3,7,10)\},\\
C_3={}&\{(0,5,8),(1,4,8),(2,6,9),(3,5,9),(3,7,11),(4,7,10)\},\\
C_4={}&\{(1,5,8),(3,6,9),(4,7,11)\}.
\end{aligned}
\]
These lists are certificate witnesses, not alternative definitions of
admissible chords or moves. Three finite propositions are evaluated by Lean's
kernel from the intrinsic definitions:
\begin{enumerate}
\item If $p\in\Sh$ and $p\to q\to r\to s$, then $q\in K$.
\item For each $C_j$, a successor of a state in $C_j$ that belongs to $K$
also belongs to $C_j$. The union of all valid marked readings of states in
$C_j$ has at most six elements.
\item For $p\to q\to r$ with $p,q,r\in K$, every choice of marked readings
at the three positions gives three pairwise distinct types.
\end{enumerate}

\begin{proof}[Proof of Theorem~\ref{thm:main}]
Project the actual walk to $\sigma(c(n))$. At its first position, use the
assumed predecessor and the next two moves in certificate 1. At any later
position, use the preceding position and the next two moves. Hence every
projected state belongs to $K$.

Choose a component list $C_j$ containing the initial state. Certificate 2
and induction keep the entire projected walk in that same list. Every marked
reading then lies in its computed union of at most six types. Certificate 3
applied to the first two edges gives three distinct types. The direct-reading
bridge transfers both conclusions to the original ordered chord walk.
\end{proof}

The component capacities, with all symmetric-root choices included, are:
\begin{center}\small
\begin{tabular}{@{}clc@{}}\toprule
List&Marked quality/inversion inventory&Capacity\\\midrule
$C_1$&$M_0,\ (+M7)_0,(+M7)_3,\ (\varnothing7)_0,(\varnothing7)_1,(\varnothing7)_3$&6\\
$C_2$&$+_0,+_1,+_2,\ (m7)_0,(m7)_3$&5\\
$C_3$&$7_0,7_2,7_3,\ m_2,\ (mM7)_0,(mM7)_3$&6\\
$C_4$&$(\circ7)_0,(\circ7)_1,(\circ7)_2,(\circ7)_3,\ (M7)_0,(M7)_3$&6\\\bottomrule
\end{tabular}
\end{center}
Subscripts are inversion indices, with zero denoting root position.

\section{Why the hypotheses are material}
Consider the actual ordered chords
\[
\begin{aligned}
&(0,0,3,6),\ (11,0,3,7),\ (10,1,3,7),\ (9,1,4,7),\\
&(8,1,4,8),\ (7,1,4,9),\ (6,1,5,9).
\end{aligned}
\]
Their moves are $S,T,A,S,S,A$. Their quality/inversion types are, in order,
\[
\mathrm{dim}_0,\ (mM7)_3,\ 7_2,\ 7_0,\ m_2,\ 7_3,\ (mM7)_0.
\]
All seven are distinct and have unambiguous roots. The final state admits an
infinite admissible continuation; the first chord has no admissible
predecessor. Thus simply discarding the initial-approachability hypothesis
would make the upper bound false even for uniquely rooted chords.

The optional formal witness supplement verifies a 36-chord actual periodic
walk attaining six types, and another attaining exactly three. It also
verifies the seven-type example above with an infinite continuation and
proves that its initial chord has no valid predecessor. Thus both endpoints
of the conditional bound are attainable.

A finite approachable walk can instead leave its recurrent family and end at
an additional type. Therefore approachability alone is insufficient if
indefinite continuation is dropped. These observations explain why our
hypotheses are stated visibly rather than inferred from the word
``progression.'' They are not by themselves counterexamples to a uniquely
specified interpretation of Pyle's conjecture.

\section{Formal implementation and verification boundary}
The core proof uses Lean 4.32.2 and its standard library. It freshly reuses
23 source modules from the earlier public project \cite{Prior}, ending in
\texttt{NodeJoin} and \texttt{LocalQuotientLift}; no earlier compiled project
objects are needed. The new modules provide the main result and concrete witnesses:
\begin{itemize}
\item \texttt{RootMarkedBounds}: template readings, finite certificates,
and the witness-plus-list-cover theorem \texttt{three\_to\_six}.
\item \texttt{RawReading}: direct actual-voice reading equivalence.\\
The export is \texttt{raw\_three\_to\_six}.
\item \texttt{Witnesses}: literal periodic chord arrays, universal periodic
path proofs, endpoint attainment, and the unapproachable seven-type example.
\item \texttt{Cardinality}: a Mathlib wrapper producing a finite set equal
to the whole infinite range, with cardinality between three and six.
\end{itemize}
The core conclusion gives three distinct attained types and a list of length
at most six covering every index; this is already the desired cardinality
bound. The wrapper merely expresses it with \texttt{Finset.card}.

The printed axioms of the three main bound exports are precisely
\begin{center}\texttt{propext}, \texttt{Classical.choice}, \texttt{Quot.sound}.\end{center}
There are no proof placeholders, custom axioms, or native-evaluation axioms
in these exports. Finite checks use \texttt{decide}. The supplied rebuild
script separates the standard-library proof from the optional Mathlib
cardinality presentation. Build success establishes the exact formal
statement above; it does not establish a priority claim or identify an
unstated musical definition with a formal predicate.

\appendix
\section{Question reserved for later correspondence}\label{app:question}
Does the seven-sonority admissible prefix in Section 4 count as a generalized
omnibus progression before its first extension? If not, which chord or
transition excludes it, or constitutes the extension?

The point is to determine whether the author's intended class excludes the
unapproachable initial chord, imposes recurrence or another condition, or
uses a different extension boundary. The theorem in this note remains valid
regardless of the answer, while its precise relationship to the original
conjecture depends on that clarification. No email or publication is part of
this deliverable.

\begin{thebibliography}{9}
\bibitem{Pyle} Hayden Pyle. Generalized Omnibus Progressions: A Mathematical
Formalization of Harmonic Progressions Exhibiting Semitonal Contrary Motion
and Tertian Sonorities. \emph{Journal MusMat} IX (2025), 27--41.
DOI: 10.46926/musmat.2025v9.27-41.
\url{https://musmat.org/journal/v9/02-Pyle.pdf}.
\bibitem{Prior} GPT-6 Astra Musical Edition project. Omnibus formalization,
public snapshot 80a1a270d9355783c6c856bc463ebc936c5ae5c5.
\url{https://github.com/NoahWLono/GPT-6-Astra-Musical-Edition}.
The reused modules provide the literal chord/move model and symbolic quotient
bridge. Their earlier graph and image-count results are not presented here as
new discoveries.
\end{thebibliography}
\end{document}
